
\documentclass{../concepts/titul_protokol_p/prepareprotokol} %všechny balíčky pro protokol + titulní strana

\usepackage{graphicx}
\usepackage{newfloat}
\usepackage{caption}
\usepackage{booktabs}
\usepackage{amssymb}
\usepackage{amsmath}
\usepackage{newfloat}
\usepackage{caption}
\usepackage{siunitx}
\usepackage{textcomp}
\usepackage{gensymb}
\sisetup{inter-unit-product =$\cdot$}
%\sisetup{separate-uncertainty=true}

\praktikum{III} %číslo praktika (I, II nebo III)
\autor{David Němec} %jméno autora
\datum{DEN.\,MĚSÍC.\,2026} %datum měření
\cislo{ČÍSLO} %číslo úlohy
\nazev{NÁZEV} %název úlohy

\DeclareFloatingEnvironment[fileext=los, placement={htbp}, name=Obrázek]{schema}

\newcommand{\nobars}{
Chybové úsečky nebyly pro~svou malou velikost pro~přehlednost vykresleny.}

\newcommand{\relative}[1]{
\frac{\sigma_{#1}}{#1}
}
\newcommand{\squared}[1]{
\left( #1 \right)^2
}
\newcommand{\relsqr}[1]{
\squared{\relative{#1}}
%\left( \frac{\sigma_{#1}}{#1} \right)^2
}

\begin{document}
\renewcommand{\figurename}{Graf}
\renewcommand{\refname}{Literatura}
\renewcommand{\today}{\number\day.\,\number\month.\,\number\year}

% add another graphics path to find ZPF.pdf for title page
% now searches for images in ./ and ../concepts/titul_protokol_p/
\graphicspath{{../concepts/titul_protokol_p/}}
\maketitle

% ========================== PRACOVNÍ ÚKOL ============================
\section{Pracovní úkol}
\begin{enumerate}
    \item
\end{enumerate}

% ============================= TEORIE ================================
\section{Teorie}
\subsection{Kapacita deskového kondenzátoru}


\subsection{Metoda měření}

\begin{schema}[htbp]
    \centering
%    \includegraphics[width=0.5\linewidth]{uloha4/fourpoints.png}
    \caption{Schéma mřížkového spektrometru \cite{pokyny}}
    \label{sch:spektrometr}
\end{schema}

\begin{figure}[htbp]
    \centering
%    \includegraphics[width=0.5\linewidth]{uloha27/schema.png}
    \caption{Jiný obrázek \cite{pokyny}}
    \label{fig:obrazek}
\end{figure}



\subsection{Měřicí přístroje a jejich nejistoty}
\begin{enumerate}
    \item PŘÍSTROJ:
    
    POZNÁMKY

\end{enumerate}


% ======================== VÝSLEDKY MĚŘENÍ ===========================
\section{Výsledky měření}


\begin{figure}
    \centering
%    \includegraphics[width=0.5\linewidth]{uloha27/GRAF1.png}
    \caption{Závislost kapacity deskového kondenzátoru na napětí.}
    \label{fig:osciloskop}
\end{figure}


% ============================ DISKUZE ===============================
\section{Diskuze}


% ============================= ZÁVĚR ================================
\section{Závěr}


% =========================== LITERATURA =============================
\begin{thebibliography}{9}

\bibitem{studijni} Kolektiv ZFP KVOF MFF UK: \textit{Měření dielektrických vlastností materiálů} [online]. 
[cit. \today], dostupné z~https://physics.mff.cuni.cz/vyuka/zfp/\_media/zadani/texty/txt\_227.pdf

\bibitem{pokyny} Kolektiv ZFP KVOF MFF UK: \textit{Měření odporu} [online]. [cit. \today], dostupné
z~https://physics.mff.cuni.cz/vyuka/zfp/\_media/zadani/texty/txt\_202.pdf

\bibitem{englich} J. Englich: \textit{Úvod do praktické fyziky I: Zpracování výsledků měření}. 1. vyd. Praha: Matfyzpress, 2006

\bibitem{mikulcak} Mikulčák a kolektiv: \textit{Matematické, fyzikální a~chemické tabulky pro~střední školy: relativní permitivita}, Praha: Státní pedagogické nakladatelství, n. p., 1988

\end{thebibliography}
\end{document}